\section{Connection to VegasCore}
\label{sec:vegascore}

A companion project, VegasCore, formalizes a core calculus over
the same event-graph carrier in Lean~4~\citep{deMouraUllrich2021,mathlib2020}.
The two artifacts are independent --- Vegas is a compiler tool;
VegasCore is a self-contained semantic construction --- but they
describe overlapping objects, and a programmer interested in
proving properties of a Vegas program will eventually want to
cross the bridge between them.  This section sketches the
correspondence, lists where the two diverge, and gives the gap
between today's state and a fully verified pipeline.

\subsection{What VegasCore is}

VegasCore is a Lean~4 library that formalizes the core
\TT{commit}/\TT{reveal}/\TT{sample}/\TT{let}/\TT{ret} calculus
underlying Vegas's source language.  Its central object is an
\TT{EventGraph} type elaborated from a checked program
(\TT{WFProgram}); the well-formedness conditions
--- visibility, freshness of binders, completeness of reveals,
normalization of sample distributions, and at-least-one-legal
action --- are statically encoded as fields of \TT{WFProgram}.
From a checked program VegasCore derives a probabilistic
\TT{Machine}, native strategy carriers, and kernel-game
presentations of the protocol; it states and proves frontier
commutation, the commit/reveal information barrier, agreement of
several outcome laws, a finite Kuhn realization theorem, and
backend-refinement transport.  All of this is mechanized in
Lean.

VegasCore's role in the project is to provide a
formally-mechanized referent for the kind of semantics that
Vegas's compiler informally implements.  Where Vegas's IR
constructor refuses to accept a graph violating the
visibility-on-reads invariant, VegasCore proves the corresponding
theorem about the elaborated graph; where Vegas's
commit-reveal-expansion pass produces an informationally
correct graph by construction, VegasCore proves frontier
commutation and the information barrier as theorems about the
abstract event graph.

\subsection{Shared vocabulary}

The two projects use the same names for the things they have in
common.  The list is short and almost exact:

\begin{center}
\begin{tabular}{ll}
\toprule
Vegas (this compiler) & VegasCore (Lean) \\
\midrule
EventGraph                  & EventGraph \\
node / nodeId               & event / node id \\
guard (\TT{scope}, \TT{expr}) & per-node guard $R$ \\
\COMMIT / \REVEAL / \PUBLIC writes & hidden + commit/reveal node kinds \\
frontier (\TT{FrontierMachine}) & frontier \\
\TT{random}                 & \TT{sample} \\
\TT{withdraw}               & \TT{Ret} \\
\bottomrule
\end{tabular}
\end{center}

For an entry-level reading of either project, this table is the
glossary: the Vegas reader who picks up the VegasCore chapter
already knows what an EventGraph is, and vice versa.

\subsection{Where they diverge, and why}

Each artifact has features the other does not.  The asymmetries
are deliberate.

\paragraph{Vegas surface keywords and macros.}  VegasCore's
surface language exposes the five core constructors directly;
Vegas's surface adds the user-facing
\TT{yield}/\TT{join}/\TT{commit}/\TT{reveal}/\TT{random}
keywords, macros for expression-level abbreviation, the
\TT{or split/burn/null} handler form, and \TT{||} subgame
transfers.  These all lower at IR construction time into the
core constructors; they exist in Vegas because they make
specifications readable, not because they enlarge the semantic
universe.

\paragraph{The commit-reveal expansion pass.}  Vegas's IR runs an
explicit rewrite (\S\ref{sec:commit-reveal}) that introduces
commit nodes for concurrent public writes.  VegasCore models
programs \emph{after} the analogous transformation: its event
graph already distinguishes commit and reveal nodes.  Vegas's
pass is a compiler operation, not a semantic notion.  Equivalence
of the pre- and post-rewrite graphs at the strategic level is
something one would want to state and prove, but as of today
that proof is not in either codebase.

\paragraph{Operational subgame layer.}  Quit handlers, bail-out
timeouts, and the resulting subgame structure are first-class in
Vegas; VegasCore presently treats quit as a strategic move only,
without modelling the deadline mechanism.  A backend-refinement
chapter in VegasCore covers a generic notion of operational
refinement that could carry the bail mechanism, but the
Vegas-to-VegasCore mapping of the timeout subgame is not yet
written down.

\paragraph{Static obligations.}  VegasCore's
\TT{WFProgram} contains four well-formedness obligations as
fields: freshness of binders, reveal completeness, normalized
distributions, and at-least-one-legal-action.  Vegas's
typechecker and IR constructor enforce the same conditions but
without exposing them under a single uniform name; they live in
the typechecker, in the EventGraph constructor, and in the
commit-reveal pass.  Lifting them under VegasCore-style names
would make the correspondence cleaner; it is on the TODO list.

\paragraph{Backends.}  Vegas emits Solidity, Vyper, Gambit,
SMT-LIB, Scribble, Bitcoin, MAID, GraphViz, and the Diamond DAG
witness records (\S\ref{sec:other:diamond}); VegasCore emits
nothing --- it is a library of definitions and theorems.  This
is the kind of asymmetry one expects: one side runs, the other
side proves.

\subsection{What is and is not proved today}

The current state of the two-sided story is the following:

\begin{itemize}\itemsep1pt
\item For programs that lower to VegasCore's core calculus, the
  semantics formalized by VegasCore is consistent with the
  EventGraph that Vegas's frontend produces \emph{by construction
  of the lowering}: the lowering uses the same dependency
  inference and visibility discipline.  The IR-level invariants
  Vegas's constructor enforces are theorems in VegasCore.
\item The various backends agree with Vegas's IR structurally:
  every backend reads the same EventGraph and lowers it
  systematically.  This is closer to a code-review property than
  to a formal one --- there is no compiler proof.
\item There is no formal proof, today, that any specific
  backend's output (Solidity, Gambit, etc.) refines the
  EventGraph in any externally meaningful sense.  The
  fellowship-funded extension of this work proposes
  Solidity-to-EventGraph refinement as the central new theorem
  on which the practical claim of ``what you analyze is what
  gets deployed'' would rest.
\item The Diamond DAG (Coq/Lean) backend
  (\S\ref{sec:other:diamond}) is a different bet: rather than
  prove correctness of the compiler once, it ships a typed
  certificate \emph{with each compiled program}.  Both routes
  --- verified compiler and verified compiled artifact --- are
  on the project's TODO list as separate workstreams.  They are
  not redundant; they cover different audiences.
\end{itemize}

In summary: Vegas can be used today as a tool
without VegasCore (the compiler stands alone), and VegasCore can
be used today as a Lean library without Vegas (the calculus
stands alone).  Connecting them with mechanically-checked
refinement proofs is the natural next step and the main subject
of ongoing work.
